\SourceLeaf{410}%
\par
Das lässt sich auch durch die folgenden sechs Gleichungen ausdrücken
\begingroup\footnotesize
\[
\begin{aligned}
\frac{\sin A}{\sin a}=\frac{\sin B}{\sin b}=\frac{\sin C}{\sin c}&=-\frac{\cos\tfrac12(A+B+C)}{2\sin\tfrac12a.\sin\tfrac12b.\sin\tfrac12c}\\
=\frac{\cos\tfrac12(-A+B+C)}{2\sin\tfrac12a.\cos\tfrac12b.\cos\tfrac12c}&=\frac{\cos\tfrac12(A-B+C)}{2\cos\tfrac12a.\sin\tfrac12b.\cos\tfrac12c}=\frac{\cos\tfrac12(A+B-C)}{2\cos\tfrac12a.\cos\tfrac12b.\sin\tfrac12c}
\end{aligned}
\]
\endgroup
oder durch diese
\[
\begin{aligned}
\sin A\sin b\sin c&=\sin B\sin a\sin c=\sin C\sin a\sin b\\
&=-4\cos\tfrac12(+A+B+C).\cos\tfrac12a.\cos\tfrac12b.\cos\tfrac12c\\
&=4\cos\tfrac12(-A+B+C).\cos\tfrac12a.\sin\tfrac12b.\sin\tfrac12c\\
&=4\cos\tfrac12(+A-B+C).\sin\tfrac12a.\cos\tfrac12b.\sin\tfrac12c\\
&=4\cos\tfrac12(+A+B-C).\sin\tfrac12a.\sin\tfrac12b.\cos\tfrac12c
\end{aligned}
\]
Diese Größe ist das Sechsfache des Volumens der Pyramide, deren Ecken die drei Eckpunkte des sphärischen Dreiecks und der Mittelpunkt der Kugel sind, wenn man den Radius der Kugel \(=1\) setzt. Ferner ist diese Größe
\[
=4\operatorname{cotang}r.\sin\tfrac12a.\sin\tfrac12b.\sin\tfrac12c
\]
wobei \(r\) der sphärische Radius des um das Dreieck beschriebenen Kreises ist. Sie ist auch
\[
=4\cos r.\sin\alpha.\sin\tfrac12b.\sin\tfrac12c
\]
wobei [\(2\alpha,2\beta,2\gamma\) die Winkel bezeichnen, die zwischen je zwei der zu den Eckpunkten \(A,B,C\) gezogenen sphärischen Radien und gegenüber den Seiten \(a,b,c\) liegen, oder] \(\alpha,\beta,\gamma\) die Winkel des ebenen Dreiecks \(ABC\) sind. Denn \(2\sin\tfrac12a,\ 2\sin\tfrac12b,\ 2\sin\tfrac12c\) sind dessen Seiten und folglich ist \(4\sin\alpha.\sin\tfrac12b.\sin\tfrac12c\) sein doppelter Flächeninhalt. Dasselbe Dreieck lässt sich zugleich als Grundfläche der obigen Pyramide mit der Höhe \(\cos r\) betrachten, woraus sich die Richtigkeit unmittelbar ergibt.
\par\Needspace{12\baselineskip}
[Aus der obigen sechsfachen Gleichung leitet \textsc{Gauss} an anderer Stelle die von ihm so oft angewandten Formeln her:]
\[
\begin{aligned}
\cos\tfrac12A.\sin\tfrac12(b-c)&=\sin\tfrac12(B-C).\sin\tfrac12a\\
\sin\tfrac12A.\sin\tfrac12(b+c)&=\cos\tfrac12(B-C).\sin\tfrac12a\\
\cos\tfrac12A.\cos\tfrac12(b-c)&=\sin\tfrac12(B+C).\cos\tfrac12a\\
\sin\tfrac12A.\cos\tfrac12(b+c)&=\cos\tfrac12(B+C).\cos\tfrac12a
\end{aligned}
\]
